\mag=1200
\documentclass{amsart}
\usepackage{amssymb,amscd}
\textwidth=130mm
\advance\hoffset by -1in
\advance\voffset by -1in

\usepackage[T2A]{fontenc}
\usepackage[utf8]{inputenc}
\usepackage[russian]{babel}

\title{Разработка цифровых двойников регламентов\\
при помощи онтологических моделей}

\author{Л.~В. Авдонина}
\address{Новосибирский национальный исследовательский государственный университет, Новосибирск}

\begin{document}

\noindent УДК 004.82

\maketitle

Цифровой двойник регламента --- динамическая исполняемая модель нормативного документа,
позволяющая компьютеру понимать правила без неоднозначностей естественного языка.
Она автоматически проверяет соответствие действий требованиям, симулирует изменения
и адаптируется к новым редакциям~\cite{palchunov}.

В работе рассматривается задача создания цифрового двойника регламента на примере
приказов о направлении студентов на практику. Разработан программный прототип
\texttt{onto-practice}, который реализует полный цикл проверки: извлечение
структурированных фактов из текста приказа с помощью больших языковых моделей,
загрузку фактов в онтологию на языке OWL~2, логический вывод с использованием
ризонера HermiT и формирование объяснимого отчёта о выявленных нарушениях.

Онтологическая модель формализует ключевые понятия предметной области (студент,
приказ, практика, руководитель, направление подготовки и др.) и связи между ними.
Для проверки документов сформулирован набор из семи правил, включающих полноту
обязательных полей, корректность дат, соблюдение минимального срока между изданием
приказа и началом практики, наличие руководителя от университета, уникальность
номеров зачётных книжек и соответствие кода направления ФГОС.

Архитектура системы спроектирована расширяемой, что позволяет в перспективе
добавлять новые типы регламентов без переписывания ядра.

Полученные результаты демонстрируют работоспособность подхода
«текст регламента $\to$ онтология $\to$ цифровой двойник регламента» и показывают
перспективы применения цифровых двойников регламентов в системах управления
учебным процессом вузов.

{\bf Ключевые слова:}
цифровой двойник, регламент, онтология, OWL~2, автоматическая проверка,
большие языковые модели.

\begin{thebibliography}{1}
\bibitem{palchunov}
Пальчунов Д.~Е., Чернявцева С.~И.
Разработка интеллектуального помощника для автоматизированного порождения
документов кафедры // Программная инженерия. 2023. Том 14, №~8. С.~388--400.
DOI: 10.17587/prin.14.388-400.
\end{thebibliography}

\medskip
\noindent Научный руководитель --- д-р физ.-мат. наук, академик РИА
Д.~Е. Пальчунов.

\end{document}
